\subsection{Productivity Effects}
\label{subsec:produktivitaetswirkung}

The evaluation compares the requirements, implemented mechanisms, and
evidence. The criteria are reuse, standardization, variability,
maintainability, testability, reproducibility, developer experience,
extensibility, and platform coverage. Complexity, operational maturity, and
documentability supplement this assessment. The project's internal
acceptance report provides project-specific evidence, but no independent
replication \parencite{kleiveist2026acceptance}.
This distinction keeps observed mechanisms, executed verification activities,
and inferred effects separate. Automation alone is not considered a measurable
productivity gain.

Here, productivity encompasses the effort required to obtain a runnable
initial project, recurring setup decisions, the development flow, and
subsequent maintenance. A single activity metric does not represent it
adequately \parencite{forsgren2021space}. Conceptually, the effect can
therefore be expressed as

\begin{center}
  \emph{Productivity benefit = avoided project effort
  $-$ template maintenance effort $-$ additional complexity}
\end{center}

In the absence of time and effort measurements, it cannot be quantified.
The relationship also shows that centralized maintenance and additional
complexity can reduce the project effort avoided.

At project initialization, \texttt{control.py} combines profile selection,
generation, system checks, and installation. The acceptance results provide
evidence for these steps across all five profiles and for the rejection of
incompatible PostgreSQL selections
\parencite{kleiveist2026acceptance}. The mixed CI results for the later HEAD
limit the transferability of this evidence. Less manual setup is therefore
plausible, but a specific amount of time saved has not been demonstrated.

During development, profile-dependent entry points standardize startup,
testing, and builds. The shared frontend, configuration contract, and clear
responsibilities encourage reuse; specialized tools and platform prerequisites
remain necessary. Central changes benefit future projects, but require shared
maintenance of profiles, generator, tests, and documentation. Projects that
have already been generated are not updated automatically.
The benefit is thus distributed across project initialization, ongoing
development, and subsequent baseline maintenance. It is particularly plausible
when several technically related projects use the same foundations.

Quantitative evidence is lacking. A future controlled evaluation would need to
compare similar projects with and without the template in terms of
initialization and onboarding time, setup errors, failed builds, proportion of
reuse, and maintenance time.
It should also record the time to the first domain-specific feature and the
number of manual setup decisions. Only such comparative data would permit a
statement about the magnitude of the effect.
\beginnerinput{workspace/statement/700_bewertung/710}
